\begin{theorem}[Theorem~1.3(iii)]
Let $\mathbb G$ be a sub-elliptic setting (\noderef{CarnotSetting}) with step $r$, frame
$\sigma$ and first-layer projection $\pi_1:\mathbb R^{n}\to\mathbb R^{m}$. Assume the cited
results of the literature \emph{at $\mathbb G$} (\noderef{CitedResults}), and also the cited
results of the literature \emph{at the Euclidean setting $\mathbb G_E$ on $\mathbb R^{m}$}
(\noderef{EuclidSetting} at $N:=m$, and \noderef{CitedResults} at that setting). Let
$A:\mathbb R^{m}\to\mathbb R^{m\times m}$ and $H:\mathbb R^{m}\to\mathbb R$ satisfy the
structure conditions (\noderef{StructureConditionsES}). Let $\Om\subset\mathbb R^{n}$ be a
bounded open set with $\mathbb R^{n}\mns\Om\ne\emptyset$, let $u,v:\mathbb R^{n}\to\mathbb R$
be continuous on $\bar\Om$ with $u$ a viscosity sub-solution (\noderef{IsViscSubsol}) and
$v$ a viscosity super-solution (\noderef{IsViscSupersol}) of $\mathcal L\cdot=0$ in $\Om$ and
$u\le v$ on $\del\Om$, and let $R\ge0$ satisfy $|u|\le R/2$ and $|v|\le R/2$ on $\bar\Om$.
Then:
\begin{enumerate}
\item if $r=1$, then $u\le v$ in $\Om$;
\item if there are $F,Gf:\mathbb R^{m}\to\mathbb R$, continuous on $\pi_1(\bar\Om)$, with
  $u=F\circ\pi_1$ and $v=Gf\circ\pi_1$, then $u\le v$ in $\Om$.
\end{enumerate}
\end{theorem}
\begin{proof}
This is Theorem~1.3(iii), both disjuncts as the paper prints them ("$r=1$, i.e. $\mathbb G$ is
the Euclidean group; more generally, $u=F\circ\pi_1$ and $v=Gf\circ\pi_1$ for continuous
$F,Gf$").

The two cited-results hypotheses are used in disjoint places, and each part uses exactly one
of them. Part (1) applies \noderef{ComparisonAbelian} to $\mathbb G$ itself and so needs the
package at $\mathbb G$. Part (2) descends to $\mathbb R^{m}$ and applies
\noderef{EuclidFirstLayerComparison}, which works inside the Euclidean setting $\mathbb G_E$
and so needs the package at $\mathbb G_E$; the package at $\mathbb G$ would not serve there,
because the three Carnot-side items of \noderef{CitedResults} are asserted only at the
setting carried as that node's parameter.

\medskip\noindent\textbf{Part (1): $r=1$.} By \noderef{CarnotSetting}, a setting of step $1$
has a commutative group law: the field $\mathtt{abelian\_of\_step\_one}$ gives
$x*y=y*x$ for all $x,y$ when $r=1$. All the remaining hypotheses of
\noderef{ComparisonAbelian} are hypotheses of the present statement --- the cited results at
$\mathbb G$, which is the setting that node is applied to here;
the structure conditions on $(A,H)$; $\Om$ bounded, open, with non-empty complement; $u,v$
continuous on $\bar\Om$; the two viscosity properties; $u\le v$ on $\del\Om$; and $R\ge0$
with $|u|,|v|\le R/2$ on $\bar\Om$. Applying \noderef{ComparisonAbelian} to
$\mathbb G$ itself therefore gives $u\le v$ in $\Om$. This is the paper's "Suppose first
$r=1$, so that $\mathbb G=(\mathbb R^{n},+)$".

\medskip\noindent\textbf{Part (2): the first-layer case.} Let $F,Gf$ be as in (2).

\emph{(2a) The degenerate corner $m=0$.} Suppose $m=0$. Then the frame is the empty family
of vector fields, so the two conditions defining an admissible family in \noderef{LieHull}
--- containing each $\sigma^{i}$, and closure under the Lie bracket --- are both vacuous for
the empty family; hence $\emptyset$ is admissible and, $\mathrm{LieHull}\ \sigma$ being the
intersection of all admissible families, $\mathrm{LieHull}\ \sigma=\emptyset$. Hörmander's
rank condition (1.6), carried as a field of \noderef{CarnotSetting}, says that for
every $x$ the span of $\{V(x):V\in\mathrm{LieHull}\ \sigma\}$ is all of $\mathbb R^{n}$; with
$\mathrm{LieHull}\ \sigma=\emptyset$ that set of values is empty and its span is $\{0\}$, so
$\mathbb R^{n}=\{0\}$ and therefore $n=0$ (for $n\ge1$ the vector $e_1$ is a non-zero element
of $\mathbb R^{n}$). So $\mathbb R^{n}$ has exactly one point, $0$; by
$\mathbb R^{n}\mns\Om\ne\emptyset$ that point is not in $\Om$, so $\Om=\emptyset$ and the
conclusion "$u\le v$ in $\Om$" holds vacuously. We may therefore assume $m\ge1$ from now on.

\emph{(2b) The descent to $\mathbb R^{m}$.} Put $\Om':=\pi_1(\Om)$. The hypotheses of
\noderef{FirstLayerEuclideanSubsol} hold for $(\mathbb G,A,H,\Om,F,Gf)$: $\Om$ is bounded and
open; $F$ and $Gf$ are continuous on $\pi_1(\bar\Om)$; $F\circ\pi_1=u$ is a viscosity
sub-solution and $Gf\circ\pi_1=v$ a viscosity super-solution of $\mathcal L\cdot=0$ in $\Om$;
and for $x\in\del\Om$, $F(\pi_1(x))=u(x)\le v(x)=Gf(\pi_1(x))$. It therefore gives
\begin{itemize}
\item[(i)] $\Om'$ is open;
\item[(ii)] $\Om'$ is bounded;
\item[(iii)] $\overline{\Om'}\subseteq\pi_1(\bar\Om)$;
\item[(v)] $F\le Gf$ on $\del\Om'$;
\item[(vi)] $F$ and $Gf$ are continuous on $\overline{\Om'}$;
\item[(vii)] the Euclidean sub-solution condition for $F$ in $\Om'$: for every $\zeta\in\Om'$,
  every $\psi$ that is $C^{2}$ on $\Om'$ with $\zeta\in\argmax_{\Om'}(F-\psi)$ and every
  second-order jet $(E,C)$ of $\psi$ at $\zeta$ (\noderef{TwiceDiffAt}) with
  $\xi(E):=(E(e_1),\ldots,E(e_m))\ne0$, one has
  $-\Tr\big(A(\xi(E))\,\mathcal C(C)\big)+H(\xi(E))\le0$, where
  $\mathcal C(C):=(C(e_i,e_j))_{i,j}$;
\item[(viii)] the dual Euclidean super-solution condition for $Gf$ in $\Om'$.
\end{itemize}
(It also gives $\del\Om'\subseteq\pi_1(\del\Om)$, which we do not need separately, (v) being
the form in which it is used.)

\emph{(2c) The complement of $\Om'$ is non-empty.} Since $m\ge1$ by (2a), the space
$\mathbb R^{m}$ is unbounded: the points $t\,e_1$, $t\in\mathbb R$, have $|t\,e_1|=|t|$,
unbounded in $t$. By (ii) the set $\Om'$ is bounded, so $\Om'\ne\mathbb R^{m}$, i.e. there is
$\zeta_0\in\mathbb R^{m}$ with $\zeta_0\notin\Om'$.

\emph{(2d) The bound on $F$ and $Gf$.} Let $\zeta\in\overline{\Om'}$. By (iii) there is
$x_0\in\bar\Om$ with $\pi_1(x_0)=\zeta$. Then $F(\zeta)=F(\pi_1(x_0))=u(x_0)$ and
$Gf(\zeta)=Gf(\pi_1(x_0))=v(x_0)$, so $|F(\zeta)|\le R/2$ and $|Gf(\zeta)|\le R/2$ by the
hypothesis $|u|,|v|\le R/2$ on $\bar\Om$.

\emph{(2e) The Euclidean comparison.} All hypotheses of
\noderef{EuclidFirstLayerComparison} are now available with the dimension $m$, the
cited-results package at the Euclidean setting $\mathbb G_E$ on $\mathbb R^{m}$ --- which is
the second cited-results hypothesis of the present statement, and is what that node requires,
since it applies \noderef{ComparisonAbelian} inside $\mathbb G_E$ and not inside
$\mathbb G$ --- the pair
$(A,H)$ (which satisfies the structure conditions by hypothesis --- note that
$\mathtt{StructureConditionsES}\ m\ A\ H$ is literally the hypothesis we carry, so no
transport is involved here), the set $\Om'$ (open by (i), bounded by (ii), with non-empty
complement by (2c)), the functions $F,Gf$ (continuous on $\overline{\Om'}$ by (vi), with
$F\le Gf$ on $\del\Om'$ by (v)), the constant $R\ge0$ with the bound (2d), and the two inline
viscosity conditions (vii) and (viii). It gives
\[
 F(\zeta)\le Gf(\zeta)\qquad\text{for every }\zeta\in\Om' .
\]
This is the paper's "By the case $r=1$ applied in $\mathbb R^{m}$, $F\le Gf$ in $\Om'$".

\emph{(2f) Descent back to $\Om$.} Let $x\in\Om$. Then $\pi_1(x)\in\pi_1(\Om)=\Om'$, so
(2e) gives $F(\pi_1(x))\le Gf(\pi_1(x))$, i.e. $u(x)\le v(x)$. As $x\in\Om$ was arbitrary,
$u\le v$ in $\Om$, which is the paper's "i.e. $u\le v$ in $\Om$".
\end{proof}
